\chapter{はじめに}

\section{研究背景}
近年,表面筋電位（surface Electromyography：sEMG）を用いたヒューマンインタフェースや義手制御,リハビリテーション支援に関する研究が盛んに行われている.sEMGは皮膚表面に電極を装着することで筋活動を非侵襲的に計測できるため,被験者への負担が少なく,日常環境での応用に適した生体信号である.

一方で,sEMG信号はノイズや電極装着状態,発汗,個人差などの影響を受けやすく,信号の再現性や安定性に課題を有している.特に単一チャネル構成のsEMGでは,取得できる情報量が限られるため,多チャネル構成と比較して分類精度の低下や環境変動への脆弱性が問題となる.このため,従来研究の多くでは複数チャネルを用いた解析や,計算資源に余裕のあるPC上での処理が主流であった.

しかし,ウェアラブルデバイスやモバイル端末の普及に伴い,エッジデバイス上でリアルタイムに動作する低コストかつ簡便なsEMG解析システムへの需要が高まっている.とりわけ,スマートフォンなどのモバイル端末上で推論を完結させるオンデバイス処理は,通信遅延やプライバシーの観点からも有効な手法である.

さらに,近年はモデルを事前に学習させるだけでなく,利用者固有のデータを用いてモデルを更新する「オンデバイス学習」に注目が集まっている.オンデバイス学習は,個人差や使用環境の変化に適応可能なモデル構築を可能とする一方で,計算資源や実装方法に制約があり,実システムへの応用事例や知見は十分に蓄積されていない.

\pagebreak

\section{研究目的}
本研究の目的は,MyoWareを用いた単一チャネルsEMGを対象として,iOS端末上でリアルタイムに動作するオンデバイス解析システムを開発するとともに,オンデバイス学習とオフラインでのファインチューニングを組み合わせた学習手法による分類精度向上の可能性を検討することである.

具体的には,Flutterを基盤としたアプリケーション構成において,BLE通信および推論処理をiOSネイティブ側で実装し,リアルタイム性と拡張性を両立したシステムを構築する.また,端末上で収集したデータを用いてモデルを更新し,その結果をオフラインで再学習する手法を導入することで,単一チャネルsEMGにおける個人差や環境変動への適応性について評価を行う.

本研究を通じて,低コストかつ実用的なオンデバイスsEMG解析システムの実現可能性を示すとともに,実装事例の少ないオンデバイス学習の応用に関する工学的知見を提供することを目指す.
